\begin{textsample}{NEG Dim 1 – vem\_ed\_19 – Score 24.00 – vem\_ed\_19/t087.txt}  \label{ex:f1_neg_146}
Entrevista Seção Conte como foi implementar GLE na DePaul e sua evolução ao longo da década.
Começamos com um programa de treinamento on-line para os docentes, com atividades predominantemente assíncronas e reuniões síncronas semanais.
Na primeira semana, os professores compartilham suas experiências; na segunda, fazem um brainstorming sobre ideias de projetos e, na terceira, abordamos o cenário intercultural, como planejar e prevenir coisas que podem dar errado ou, se estivermos em uma situação difícil, como gerenciar a experiência dos alunos e do professor.
Desde o início, fizemos parceria com o Center for Teaching and Learning, unidade da DePaul University que apoia professores com o design instrucional para cursos on-line, oferecidos desde muito antes da pandemia.
Essa equipe ajuda com aspectos tecnológicos e treinamento em tópicos de ensino-aprendizagem úteis para os professores.
O mais recente, por exemplo, foi sobre Inteligência Artificial.
Entre as diversas atividades do nosso Global Engagement Office está o treinamento de professores, o que é muito importante, especialmente para aqueles que estão começando.
Por muitos anos, nosso foco principal foi nos professores, pois sem eles não tocamos o programa.
Mais recentemente, comecei a me preocupar com a forma de chegar aos estudantes.
Muitos deles não sabem que vão participar de um Intercâmbio Virtual até o professor contar.
Então, estamos trabalhando em um pequeno módulo de preparação com informações básicas sobre o que é COIL / GLE.
Qual foi o pedido de parceria mais desafiador?
Frequentemente, os projetos mais desafiadores são aqueles com disciplinas muito específicas, como aconteceu com um professor de Matemática — seguimos buscando parceria.
Quando os docentes são novos na abordagem, há muito para produzir e planejar.
Depois que ganham experiência, percebem os benefícios de projetos interdisciplinares.
Seguimos aprendendo sobre o difícil processo de encontrar parceiros (matchmaking).
Há muitas coisas a considerar e que precisam ser resolvidas.
Uma delas é pedir informações específicas: breve descrição da disciplina, plano de ensino, tamanho típico da turma e qual o semestre ou período no qual é oferecida.

% matched lemmas: 
\end{textsample}
